\section*{第 \ref{ch:g7:motion-average-speed} 章 --- 运动图像与平均速度}

\begin{solution}{exo:g7:motion-average-speed:1}
$v = d/t$；单位 \unit{km/h} 和 \unit{m/s}。距离：
$d = v \times t$；时间：$t = d/v$。
\end{solution}

\begin{solution}{exo:g7:motion-average-speed:2}
渡船：$45 \div 3 = \qty{15}{km/h}$。野兔：$100 \div 8 =
\qty{12.5}{m/s}$。
\end{solution}

\begin{solution}{exo:g7:motion-average-speed:3}
$5 \times 3.6 = \qty{18}{km/h}$; \ $72 \div 3.6 = \qty{20}{m/s}$.
\end{solution}

\begin{solution}{exo:g7:motion-average-speed:4}
直线表示匀速运动；越陡表示越快；水平段表示静止。一
段逐渐变平的弯曲表示正在减速。
\end{solution}

\begin{solution}{exo:g7:motion-average-speed:5}
从 $35$ 分钟到 $60$ 分钟 --- 第三段：
在上升的各段中它最平缓，
直接从较小的斜率即可看出。
\end{solution}

\begin{solution}{exo:g7:motion-average-speed:6}
匀速 --- 一条直线。它每半小时
增加 \qty{2}{km}，于是 $v = \qty{4}{km/h}$。
\end{solution}

\begin{solution}{exo:g7:motion-average-speed:7}
总路程除以总时间（包括停留）。全程：
$12 \div 4 = \qty{3}{km/h}$。仅算步行：$12 \div 3 =
\qty{4}{km/h}$。
\end{solution}

\begin{solution}{exo:g7:motion-average-speed:8}
路程相等时，慢速段只是持续时间更长 ---
行程的时间中有更多花在慢速上，所以
慢速在整个行程的结果中占更大的分量。
\end{solution}

\begin{solution}{exo:g7:motion-average-speed:9}
中点猜测：\qty{22.5}{km/h}。实际计算：去程 $30 \div 15 = 2$
小时，回程 $30 \div 30 = 1$ 小时；总共 \qty{60}{km} 用时 $3$
小时，平均为 \qty{20}{km/h}。慢速的那一半占了
三小时中的两个小时。
\end{solution}

\begin{solution}{exo:g7:motion-average-speed:10}
图像：沿直线到 (\qty{30}{min}, \qty{20}{km})；水平段持续到
$45$ 分钟；再沿较陡的直线到 (\qty{60}{min}, \qty{35}{km})。
平均：\qty{35}{km} 用时一小时 --- \qty{35}{km/h}。
\end{solution}

\begin{solution}{exo:g7:motion-average-speed:11}
$d = 340 \times 6 = \qty{2040}{m}$ --- 约 \qty{2}{km}。
旧的经验法则：$6 \div 3 = 2$ 千米 --- 两者一致，
因为 \qty{340}{m/s} 几乎等于每三秒
走一千米。
\end{solution}

\begin{solution}{exo:g7:motion-average-speed:12}
前半程：\qty{30}{km}，速度 \qty{30}{km/h}，恰好需要 $1$
小时。全程 \qty{60}{km} 平均 \qty{60}{km/h}，总时间恰好只允许
为 $1$ 小时——时间已耗尽，一秒不剩。
余下的 \qty{30}{km} 需要用时 $0$ 小时：无论
多快的速度，都无法挽救平均速度；时间预算已经用完。
\end{solution}

\begin{solution}{pb:g7:motion-average-speed:1}
\textbf{1.} 快速骑行（$30$ 分钟）、$15$ 分钟停留
（送货）、较慢骑行（$30$ 分钟）、第二次 $15$ 分钟
停留 --- 两段骑行，两次送货。
\textbf{2.} 半小时行驶 \qty{9}{km}：\qty{18}{km/h}。
\textbf{3.} $15 - 9 = \qty{6}{km}$，半小时内：
\qty{12}{km/h}。
\textbf{4.} 行驶 \qty{15}{km}，用时 $1.5$ 小时：\qty{10}{km/h}。
\textbf{5.} $t = 12 \div 24 = 0.5$ 小时，即 $30$ 分钟。
\textbf{6.} $20$ 分钟是一小时的三分之一：$d = 18 \times
\tfrac{1}{3} = \qty{6}{km}$。
\textbf{7.} $25$ 分钟即 $\tfrac{25}{60}$ 小时；$v = 10
\div \tfrac{25}{60} = \qty{24}{km/h}$。
\textbf{8.} 路程：$12 + 6 + 10 = \qty{28}{km}$；时间：
$30 + 20 + 25 = 75$ 分钟 $= 1.25$ 小时；滚动平均速度：
$28 \div 1.25 = \qty{22.4}{km/h}$。
\textbf{9.} $37 \div 3 \approx \qty{12.3}{km/h}$ --- 远低于
$20$：不发奖金（她的停靠是正经工作，但规则把
停靠时间也算了进去）。
\textbf{10.} 马克斯：快速段 $12 \div 30 = 0.4$ 小时；慢速段
$12 \div 15 = 0.8$ 小时；总共 \qty{24}{km} 用时 $1.2$ 小时：恰好
\qty{20}{km/h} --- 并非\emph{高于} $20$：不发奖金。
\textbf{11.} 他的 $22.5$ 是对两个\emph{数值}取了平均，仿佛
每种速度各占一半时间。实际上慢速段占去了
$0.8$ 小时，即他 $1.2$ 小时的三分之二，
在行程的分量中同样占了三分之二，并把真实平均值拉低到 $20$。
\textbf{12.} 例如：奖金依\emph{滚动}
平均速度发放 --- 即总路程仅除以骑行时间，
图像中的水平段不计入。它衡量的是
车轮实际保持的速度，奖励快速骑行，
而不是跳过送货。
\textbf{13.} “直线段表示稳定的速度；水平段表示
静止但时钟仍在走；较陡的段表示
更快的速度 --- 陡峭程度即速度。切勿
对平均速度的数值再取平均：应对整段行程取平均 --- 总路程
除以总时间 --- 否则慢速时段会让
这张表格出洋相。”
\end{solution}
